\section{Related Work}
\label{sec:related}
This section relates AgentWorkShop to plant-facing industrial agents and to
production and automation abstractions for context, tool access, and
operational constraints. Table~\ref{tab:nearest} compares three systems
selected to represent industrial control generation, factory-module agent
coordination, and supervisory service reconfiguration. It uses execution
scope, stated integration boundary, and evaluation focus as common dimensions.
This is a focused comparison rather than an exhaustive survey, and it does
not treat an unreported feature as absent.

\subsection{Industrial Agents and Plant Integration}
Industrial-agent systems span end-to-end automation control, supervisory
services, programmable logic controller (PLC) interfaces, manufacturing
coordination, and operator-facing machine access. Xia et al. present an
LLM-agent framework for end-to-end industrial automation control, using
structured prompts and event-driven context for production planning and
operation \cite{xia2025control}. Huang et al. present a multimodal
architecture whose stated focus is automated control generation for
industrial cyber-physical systems \cite{huang2026made}. Biskupovic et al.
study agent-based supervision for service-oriented industrial systems
\cite{biskupovic2025supervision};
Agents4PLC and InstructMPC address PLC code generation and human--LLM
context-aware control, respectively
\cite{liu2026agents4plc,wu2025instructmpc}. Manufacturing studies consider
dynamic resource-capability exploration \cite{lim2025llmenabled}, multi-agent
shop-floor coordination \cite{zhao2026shopfloor}, and digital-twin-supported
fault handling \cite{gill2025llmdt}. Natural-language machine access
\cite{hofmann2025nlc} and real-time LLM control of a robot-operated serial
line \cite{waseem2026pretrained} extend this work to operator interfaces and
robotic-line control. Together, these studies address different operational
contexts, interfaces, and allocations of responsibility between agents,
operators, and controllers.
\begin{table*}[!t]
\centering
\caption{Selected plant-facing agent systems and their integration settings.
The table compares the stated execution scope and evaluation focus; it is not
a performance ranking and does not infer unreported features.}
\label{tab:nearest}
\small
\setlength{\tabcolsep}{4pt}
\begin{tabularx}{\textwidth}{@{}p{2.3cm}p{2.5cm}XX@{}}
\toprule
System & Execution unit & Stated integration boundary & Evaluation focus \\
\midrule
Huang et al. \cite{huang2026made} & Automated industrial control generation & Multimodal industrial CPS architecture; the paper's stated focus is generating control actions & Control generation is its stated scope; no head-to-head comparison is made here \\
K{\"o}hle et al. \cite{koehle2026llmdriven} & Factory modules; orchestrator, peer-to-peer, and monolithic designs & Per-module MCP servers expose OPC UA methods, with MQTT coordination, injected factory state, and state-based tool filtering & Nine simulated production challenges, including silent-fault detection \\
Biskupovic et al. \cite{biskupovic2025supervision} & Service selection and reconfiguration in service-oriented ICPSs & Predefined component library; performance and computational-load criteria; online monitoring & Reselection after poor performance; reported adaptability and performance results \\
\midrule
AgentWorkShop (this work) & Collaboration task and registered harness call & Line--product--recipe--run context; per-call policy checks; semantic process nodes resolve to configured protocol drivers & Two configured simulator scenes exercise the task-to-device path; command receipts and process observations remain distinct, context-linked records \\
\bottomrule
\end{tabularx}
\end{table*}

These systems target different levels of the industrial-agent stack: Huang et
al. state a control-generation focus, K{\"o}hle et al. coordinate agents around
factory modules, and Biskupovic et al. reconfigure service compositions
\cite{huang2026made,koehle2026llmdriven,biskupovic2025supervision}. This paper
examines the surrounding integration path: line--product--recipe--run binding,
per-call authorization, semantic node-to-driver routing, and distinct command
and process records. These are scope dimensions, not claims that prior work
lacks any one mechanism. With no head-to-head experiment, the contribution is
the implemented composition and evidence boundary, not comparative advantage
or a new reasoning/control algorithm. Agent strategy remains replaceable.
Authority and command confirmation are operational services consistent with
safety-aware decision support \cite{jiang2026agentic}; they do not replace
process protection or establish operator effectiveness.

\subsection{Production Context, Agent Interfaces, and Automation}
Industrial integration also depends on abstractions that are not specific to
agent architectures. MCP standardizes tool discovery and calls with
request-scoped authorization \cite{mcp2024}, whereas
ISA-88 and ISA-95 describe recipes and production operations
\cite{isa88_2010,isa95_2025}. Run-to-run methods connect process observations
to successive production runs \cite{moyne2000r2r}. At the agent level,
AgentSpec and CaMeL address tool behavior and capability-based information
flow \cite{wang2025agentspec,debenedetti2025camel}. Together, these works
provide complementary abstractions for production context, tool interaction,
and information flow. They do not, by themselves, specify the shared
integration architecture presented here.

Industrial security and record requirements further constrain that
architecture: FDA guidance discusses trustworthy records and audit trails
\cite{part11}, IEC~62443-3-3 specifies system-security requirements
\cite{iec62443-3-3}, and DNP3 includes select-before-operate
\cite{ieee1815dnp3}. Control-assurance methods address state evolution through
shielding, control-barrier functions, or Simplex architectures
\cite{alshiekh2018shielding,ames2017cbf,sha2001simplex}. These standards and
methods inform the boundaries within which an industrial-agent framework must
operate; they are not contributions claimed by AgentWorkShop.

The evaluation tests selected software paths and simulator scenes, not
physical-plant performance.
